% Options for packages loaded elsewhere
\PassOptionsToPackage{unicode}{hyperref}
\PassOptionsToPackage{hyphens}{url}
%
\documentclass[
  doc]{apa6}
\usepackage{amsmath,amssymb}
\usepackage{iftex}
\ifPDFTeX
  \usepackage[T1]{fontenc}
  \usepackage[utf8]{inputenc}
  \usepackage{textcomp} % provide euro and other symbols
\else % if luatex or xetex
  \usepackage{unicode-math} % this also loads fontspec
  \defaultfontfeatures{Scale=MatchLowercase}
  \defaultfontfeatures[\rmfamily]{Ligatures=TeX,Scale=1}
\fi
\usepackage{lmodern}
\ifPDFTeX\else
  % xetex/luatex font selection
\fi
% Use upquote if available, for straight quotes in verbatim environments
\IfFileExists{upquote.sty}{\usepackage{upquote}}{}
\IfFileExists{microtype.sty}{% use microtype if available
  \usepackage[]{microtype}
  \UseMicrotypeSet[protrusion]{basicmath} % disable protrusion for tt fonts
}{}
\makeatletter
\@ifundefined{KOMAClassName}{% if non-KOMA class
  \IfFileExists{parskip.sty}{%
    \usepackage{parskip}
  }{% else
    \setlength{\parindent}{0pt}
    \setlength{\parskip}{6pt plus 2pt minus 1pt}}
}{% if KOMA class
  \KOMAoptions{parskip=half}}
\makeatother
\usepackage{xcolor}
\usepackage{graphicx}
\makeatletter
\def\maxwidth{\ifdim\Gin@nat@width>\linewidth\linewidth\else\Gin@nat@width\fi}
\def\maxheight{\ifdim\Gin@nat@height>\textheight\textheight\else\Gin@nat@height\fi}
\makeatother
% Scale images if necessary, so that they will not overflow the page
% margins by default, and it is still possible to overwrite the defaults
% using explicit options in \includegraphics[width, height, ...]{}
\setkeys{Gin}{width=\maxwidth,height=\maxheight,keepaspectratio}
% Set default figure placement to htbp
\makeatletter
\def\fps@figure{htbp}
\makeatother
\setlength{\emergencystretch}{3em} % prevent overfull lines
\providecommand{\tightlist}{%
  \setlength{\itemsep}{0pt}\setlength{\parskip}{0pt}}
\setcounter{secnumdepth}{-\maxdimen} % remove section numbering
% Make \paragraph and \subparagraph free-standing
\ifx\paragraph\undefined\else
  \let\oldparagraph\paragraph
  \renewcommand{\paragraph}[1]{\oldparagraph{#1}\mbox{}}
\fi
\ifx\subparagraph\undefined\else
  \let\oldsubparagraph\subparagraph
  \renewcommand{\subparagraph}[1]{\oldsubparagraph{#1}\mbox{}}
\fi
\newlength{\cslhangindent}
\setlength{\cslhangindent}{1.5em}
\newlength{\csllabelwidth}
\setlength{\csllabelwidth}{3em}
\newlength{\cslentryspacingunit} % times entry-spacing
\setlength{\cslentryspacingunit}{\parskip}
\newenvironment{CSLReferences}[2] % #1 hanging-ident, #2 entry spacing
 {% don't indent paragraphs
  \setlength{\parindent}{0pt}
  % turn on hanging indent if param 1 is 1
  \ifodd #1
  \let\oldpar\par
  \def\par{\hangindent=\cslhangindent\oldpar}
  \fi
  % set entry spacing
  \setlength{\parskip}{#2\cslentryspacingunit}
 }%
 {}
\usepackage{calc}
\newcommand{\CSLBlock}[1]{#1\hfill\break}
\newcommand{\CSLLeftMargin}[1]{\parbox[t]{\csllabelwidth}{#1}}
\newcommand{\CSLRightInline}[1]{\parbox[t]{\linewidth - \csllabelwidth}{#1}\break}
\newcommand{\CSLIndent}[1]{\hspace{\cslhangindent}#1}
\ifLuaTeX
\usepackage[bidi=basic]{babel}
\else
\usepackage[bidi=default]{babel}
\fi
\babelprovide[main,import]{english}
% get rid of language-specific shorthands (see #6817):
\let\LanguageShortHands\languageshorthands
\def\languageshorthands#1{}
% Manuscript styling
\usepackage{upgreek}
\captionsetup{font=singlespacing,justification=justified}

% Table formatting
\usepackage{longtable}
\usepackage{lscape}
% \usepackage[counterclockwise]{rotating}   % Landscape page setup for large tables
\usepackage{multirow}		% Table styling
\usepackage{tabularx}		% Control Column width
\usepackage[flushleft]{threeparttable}	% Allows for three part tables with a specified notes section
\usepackage{threeparttablex}            % Lets threeparttable work with longtable

% Create new environments so endfloat can handle them
% \newenvironment{ltable}
%   {\begin{landscape}\centering\begin{threeparttable}}
%   {\end{threeparttable}\end{landscape}}
\newenvironment{lltable}{\begin{landscape}\centering\begin{ThreePartTable}}{\end{ThreePartTable}\end{landscape}}

% Enables adjusting longtable caption width to table width
% Solution found at http://golatex.de/longtable-mit-caption-so-breit-wie-die-tabelle-t15767.html
\makeatletter
\newcommand\LastLTentrywidth{1em}
\newlength\longtablewidth
\setlength{\longtablewidth}{1in}
\newcommand{\getlongtablewidth}{\begingroup \ifcsname LT@\roman{LT@tables}\endcsname \global\longtablewidth=0pt \renewcommand{\LT@entry}[2]{\global\advance\longtablewidth by ##2\relax\gdef\LastLTentrywidth{##2}}\@nameuse{LT@\roman{LT@tables}} \fi \endgroup}

% \setlength{\parindent}{0.5in}
% \setlength{\parskip}{0pt plus 0pt minus 0pt}

% Overwrite redefinition of paragraph and subparagraph by the default LaTeX template
% See https://github.com/crsh/papaja/issues/292
\makeatletter
\renewcommand{\paragraph}{\@startsection{paragraph}{4}{\parindent}%
  {0\baselineskip \@plus 0.2ex \@minus 0.2ex}%
  {-1em}%
  {\normalfont\normalsize\bfseries\itshape\typesectitle}}

\renewcommand{\subparagraph}[1]{\@startsection{subparagraph}{5}{1em}%
  {0\baselineskip \@plus 0.2ex \@minus 0.2ex}%
  {-\z@\relax}%
  {\normalfont\normalsize\itshape\hspace{\parindent}{#1}\textit{\addperi}}{\relax}}
\makeatother

\makeatletter
\usepackage{etoolbox}
\patchcmd{\maketitle}
  {\section{\normalfont\normalsize\abstractname}}
  {\section*{\normalfont\normalsize\abstractname}}
  {}{\typeout{Failed to patch abstract.}}
\patchcmd{\maketitle}
  {\section{\protect\normalfont{\@title}}}
  {\section*{\protect\normalfont{\@title}}}
  {}{\typeout{Failed to patch title.}}
\makeatother

\usepackage{xpatch}
\makeatletter
\xapptocmd\appendix
  {\xapptocmd\section
    {\addcontentsline{toc}{section}{\appendixname\ifoneappendix\else~\theappendix\fi: #1}}
    {}{\InnerPatchFailed}%
  }
{}{\PatchFailed}
\makeatother
\usepackage{csquotes}
\ifLuaTeX
  \usepackage{selnolig}  % disable illegal ligatures
\fi
\IfFileExists{bookmark.sty}{\usepackage{bookmark}}{\usepackage{hyperref}}
\IfFileExists{xurl.sty}{\usepackage{xurl}}{} % add URL line breaks if available
\urlstyle{same}
\hypersetup{
  pdfauthor={Frederik Aust1,9, Tobias Heycke1,2, Benedek Kurdi3,10,12, Pieter Van Dessel5, Jeremy Cone8, Melissa J. Ferguson10, Xiaoqing Hu6, Congjiao Jiang4,14, Robert J. Rydell7, Lisa Spitzer1,13, Christoph Stahl1, Christine Vitiello4,11, \& Jan De Houwer5},
  pdflang={en-EN},
  hidelinks,
  pdfcreator={LaTeX via pandoc}}

\title{Supplementary online material for

\emph{Of Two Minds: A registered replication}}
\author{Frederik Aust\textsuperscript{1,9}, Tobias Heycke\textsuperscript{1,2}, Benedek Kurdi\textsuperscript{3,10,12}, Pieter Van Dessel\textsuperscript{5}, Jeremy Cone\textsuperscript{8}, Melissa J. Ferguson\textsuperscript{10}, Xiaoqing Hu\textsuperscript{6}, Congjiao Jiang\textsuperscript{4,14}, Robert J. Rydell\textsuperscript{7}, Lisa Spitzer\textsuperscript{1,13}, Christoph Stahl\textsuperscript{1}, Christine Vitiello\textsuperscript{4,11}, \& Jan De Houwer\textsuperscript{5}}
\date{}


\shorttitle{Replication of Rydell et al.}

\authornote{

Correspondence concerning this article should be addressed to Frederik Aust, Herbert-Lewin-Str. 2, 50931 Cologne, Germany. E-mail: \href{mailto:frederik.aust@uni-koeln.de}{\nolinkurl{frederik.aust@uni-koeln.de}}

}

\affiliation{\vspace{0.5cm}\textsuperscript{1} University of Cologne\\\textsuperscript{2} Leibniz Institute for the Social Sciences (GESIS)\\\textsuperscript{3} Harvard University\\\textsuperscript{4} University of Florida\\\textsuperscript{5} Ghent University\\\textsuperscript{6} The University of Hong Kong\\\textsuperscript{7} Indiana University\\\textsuperscript{8} Williams College\\\textsuperscript{9} University of Amsterdam\\\textsuperscript{10} Yale University\\\textsuperscript{11} Air Force Research Laboratory\\\textsuperscript{12} University of Illinois Urbana--Champaign\\\textsuperscript{13} Leibniz Institute for Psychology (ZPID)\\\textsuperscript{14} Peking University}

\begin{document}
\maketitle

\newpage

In the following we report details on our power analyses, the model specifications used for the Bayesian model comparisons, and additional secondary analyses.

\hypertarget{sensitivity-of-our-design}{%
\section{Sensitivity of our design}\label{sensitivity-of-our-design}}

As noted in the main text, our assessment of the statistical senstivitiy of our design focused the tests of simple \emph{learning block} effects, because they are of primary theoretical interest and less sensitive than the test of the interaction.
In Experiment 1, our planned contrasts had 95\% power to detect learning block effects as small as \(\delta_z = 0.42\) (\(\eta_p^2 = .081\)).
In Experiment 2, our planned contrasts had 95\% power to detect learning block effects as small as \(\delta_z = 0.40\) (\(\eta_p^2 = .040\)) or as small as \(\delta_z = 0.29\) (\(\eta_p^2 = .020\)) and \(\delta_z = 0.20\) (\(\eta_p^2 = .010\)) when pooling participants across one or both between-participant factors (\(N = 1280\), \(\alpha = .05\), two-sided tests).
Figure~\ref{fig:power-curve-plot} illustrates the implied sensitivity of our designs in units of Cohen's \(\delta\) as a function of the assumed repeated-measures correlation \(\rho\).
Note that these are conservative estimates as they do not take into account the additional 60 participants from our pilot study that we included in the analysis.








\begin{figure}
\centering
\includegraphics{supplements_files/figure-latex/power-curve-plot-1.pdf}
\caption{\label{fig:power-curve-plot}Sensitivity curves for simple learning block effects in Experiment 1 (\textbf{A}) and Experiment 2 (\textbf{B}) as a function of the assumed repeated measures correlation.
Simple contrast are estimated within each \emph{flashed word presentation duration}-\emph{valence order} combination (not adjusted for multiple comparisons).
Collapsed contrasts pool participants across one or both of the between subject factors (i.e., \emph{flashed word presentation duration} and \emph{valence order}).
\(\delta_z\)-values in the legend correspond to 95\% and 80\% power.
Note that these are conservative estimates as they do not take into account the additional 60 participants from our pilot study that we included in the analysis.
The grey dashed line at the top represents the smallest learning block difference reported by Rydell et al.~{[}\(d_z = 0.47\); Rydell, McConnell, Mackie, and Strain (2006){]}.}
\end{figure}

\hypertarget{experiment-1}{%
\section{Experiment 1}\label{experiment-1}}

In the following we report details on our power analysis, the model specification used for the Bayesian model comparisons, and additional secondary analyses.
We report results from the linear mixed model analysis of the IAT response times, from prior sensitivity analyses for the Bayesian model comparisons, and from an exploratory analysis of the relationship between the recognition accuracy of briefly flashed words and indirectyl measured evaluations.
Table~\ref{tab:otm1-participant-table} summarizes the participants' demographics separately for each location of data collection.





\begin{table}[h]

\begin{center}
\begin{threeparttable}

\caption{\label{tab:otm1-participant-table}Participant demographics by location.}

\begin{tabular}{lrrrrr}
\toprule
 &  &  & \multicolumn{3}{c}{Gender} \\
\cmidrule(r){4-6}
Location & \multicolumn{1}{c}{$n$} & \multicolumn{1}{c}{Age} & \multicolumn{1}{c}{Female (\%)} & \multicolumn{1}{c}{Male (\%)} & \multicolumn{1}{c}{Non-binary (\%)}\\
\midrule
Cologne & 51 & 24.61 [18, 64] & 70.59 & 27.45 & 1.96\\
Ghent & 50 & 21.92 [17, 50] & 84.00 & 16.00 & 0.00\\
Harvard & 52 & 19.58 [18, 22] & 59.62 & 38.46 & 1.92\\
\bottomrule
\addlinespace
\end{tabular}

\begin{tablenotes}[para]
\normalsize{\textit{Note.} Mean age is given with range in brackets.}
\end{tablenotes}

\end{threeparttable}
\end{center}

\end{table}

\hypertarget{mixed-model-analysis}{%
\subsection{Mixed model analysis}\label{mixed-model-analysis}}

The ANOVA of IAT scores reported in the main text ignores potential systematic trial-to-trial variability in IAT response latencies due to stimuli.
Any such systematic but unaccounted-for variance can inflate test statistics and yield anticonservative \(p\) values as well as underestimated confidence intervals.
We, therefore, also conducted a linear mixed model analysis of response times with crossed random effects for participants and items to ensure that our conclusion are not contingent on inadvertent stimulus effects (for details see Wolsiefer, Westfall, \& Judd, 2017).
For this analysis we excluded participants with error rates across all blocks larger than 50\% or who responded faster than 300 ms on at least 10\% of all trials.
We additionally discarded trials in which responses were faster than 400 ms or slower than 10 s.
These exclusion criteria are the same as those used by Wolsiefer et al. (2017).

We analyzed standardized response latencies, that is, the time that elapsed between stimulus presentation and \emph{correct} response divided by the standard deviation of all response latencies in a given block, Figure~\ref{fig:otm1-iat-plot}.
To assess the reversal of the response mapping effect, we contrasted the common response mapping of Bob and negative words with the common mapping of Bob and positive words.
Hence, larger values indicate more favorable evaluations.




\begin{figure}
\centering
\includegraphics{supplements_files/figure-latex/otm1-iat-plot-1.pdf}
\caption{\label{fig:otm1-iat-plot}Standardized IAT response latencies across learning blocks.
Black-rimmed points represent condition means, error bars represent 95\% bootstrap confidence intervals based on 10,000 samples.}
\end{figure}





\begin{lltable}

\begin{TableNotes}[para]
\normalsize{\textit{Note.} The model additionally included random participant and item effects with random intercepts and random slopes for all manipulations during the learning procedure and their interactions.}
\end{TableNotes}

\begin{longtable}{lrrrrr}\noalign{\getlongtablewidth\global\LTcapwidth=\longtablewidth}
\caption{\label{tab:otm1-lmer-table-fixed-effects}Fixed effect estimates of the linear mixed model analysis of standardized IAT response times.}\\
\toprule
Term & \multicolumn{1}{c}{$\hat{\beta}$} & \multicolumn{1}{c}{95\% CI} & \multicolumn{1}{c}{$t$} & \multicolumn{1}{c}{$\mathit{df}$} & \multicolumn{1}{c}{$p$}\\
\midrule
\endfirsthead
\caption*{\normalfont{Table \ref{tab:otm1-lmer-table-fixed-effects} continued}}\\
\toprule
Term & \multicolumn{1}{c}{$\hat{\beta}$} & \multicolumn{1}{c}{95\% CI} & \multicolumn{1}{c}{$t$} & \multicolumn{1}{c}{$\mathit{df}$} & \multicolumn{1}{c}{$p$}\\
\midrule
\endhead
Intercept & 2.11 & {}[1.89, 2.33] & 18.69 & 96.14 & < .001\\
Response mapping & -0.08 & {}[-0.20, 0.04] & -1.37 & 17.28 & .188\\
Learning block & 0.36 & {}[0.14, 0.58] & 3.21 & 114.15 & .002\\
Valence order & 0.12 & {}[-0.17, 0.41] & 0.83 & 96.45 & .407\\
Category & 0.32 & {}[0.17, 0.48] & 4.13 & 24.54 & < .001\\
Word type & 0.06 & {}[-0.01, 0.12] & 1.69 & 84.92 & .094\\
Image type & -0.15 & {}[-0.20, -0.10] & -5.79 & 18539.18 & < .001\\ \midrule
Response mapping $\times$ Learning block & -0.18 & {}[-0.32, -0.04] & -2.59 & 15.97 & .020\\
Response mapping $\times$ Valence order & -0.28 & {}[-0.45, -0.11] & -3.16 & 19.55 & .005\\
Learning block $\times$ Valence order & 0.01 & {}[-0.32, 0.35] & 0.08 & 52.85 & .940\\
Response mapping $\times$ Category & -0.02 & {}[-0.14, 0.10] & -0.33 & 15.48 & .748\\
Response mapping $\times$ Word type & -0.03 & {}[-0.10, 0.04] & -0.76 & 229.26 & .451\\
Response mapping $\times$ Image type & 0.06 & {}[-0.01, 0.13] & 1.69 & 16907.60 & .091\\
Learning block $\times$ Category & 0.09 & {}[-0.02, 0.20] & 1.61 & 15.18 & .129\\
Learning block $\times$ Word type & 0.00 & {}[-0.07, 0.07] & 0.09 & 289.39 & .931\\
Learning block $\times$ Image type & -0.01 & {}[-0.09, 0.06] & -0.40 & 18911.06 & .689\\
Valence order $\times$ Category & 0.08 & {}[-0.10, 0.25] & 0.86 & 16.71 & .403\\
Valence order $\times$ Word type & 0.00 & {}[-0.08, 0.08] & -0.05 & 89.72 & .957\\
Valence order $\times$ Image type & 0.05 & {}[-0.02, 0.12] & 1.37 & 19279.71 & .171\\ \midrule
Response mapping $\times$ Learning block $\times$ Valence order & 0.33 & {}[0.11, 0.54] & 2.99 & 8.40 & .016\\
Response mapping $\times$ Learning block $\times$ Category & -0.12 & {}[-0.28, 0.04] & -1.43 & 26.09 & .165\\
Response mapping $\times$ Learning block $\times$ Word type & 0.02 & {}[-0.08, 0.12] & 0.31 & 442.51 & .759\\
Response mapping $\times$ Learning block $\times$ Image type & 0.02 & {}[-0.08, 0.13] & 0.48 & 20852.22 & .634\\
Response mapping $\times$ Valence order $\times$ Category & -0.08 & {}[-0.26, 0.10] & -0.89 & 17.69 & .383\\
Response mapping $\times$ Valence order $\times$ Word type & 0.00 & {}[-0.10, 0.10] & 0.01 & 228.77 & .990\\
Response mapping $\times$ Valence order $\times$ Image type & 0.00 & {}[-0.11, 0.10] & -0.06 & 15855.48 & .954\\
Learning block $\times$ Valence order $\times$ Category & -0.17 & {}[-0.37, 0.04] & -1.60 & 7.83 & .150\\
Learning block $\times$ Valence order $\times$ Word type & -0.03 & {}[-0.14, 0.08] & -0.58 & 66.23 & .561\\
Learning block $\times$ Valence order $\times$ Image type & -0.04 & {}[-0.14, 0.07] & -0.68 & 14359.59 & .498\\ \midrule
Response mapping $\times$ Learning block $\times$ Valence order $\times$ Category & 0.25 & {}[0.00, 0.50] & 1.98 & 12.50 & .071\\
Response mapping $\times$ Learning block $\times$ Valence order $\times$ Word type & 0.06 & {}[-0.08, 0.21] & 0.85 & 166.33 & .397\\
Response mapping $\times$ Learning block $\times$ Valence order $\times$ Image type & -0.01 & {}[-0.15, 0.14] & -0.07 & 14550.19 & .945\\
\bottomrule
\addlinespace
\insertTableNotes
\end{longtable}

\end{lltable}





\begin{lltable}

\begin{TableNotes}[para]
\normalsize{\textit{Note.} We report the estimated standard deviations in the main diagonals and the correlations in the off-diagnoals. The percentages of variance for the random effects were calculated by dividing each variance component by the total random variance, i.e., the sum of the random-effect variances.}
\end{TableNotes}

\small{

\begin{longtable}{lcrrrrrrrr}\noalign{\getlongtablewidth\global\LTcapwidth=\longtablewidth}
\caption{\label{tab:otm1-lmer-table-random-effects}Random effect estimates and correlations of the linear mixed model analysis of standardized IAT response times.}\\
\toprule
 & \multicolumn{1}{c}{\% of variance} & \multicolumn{1}{c}{1. } & \multicolumn{1}{c}{2. } & \multicolumn{1}{c}{3. } & \multicolumn{1}{c}{4. } & \multicolumn{1}{c}{5. } & \multicolumn{1}{c}{6. } & \multicolumn{1}{c}{7. } & \multicolumn{1}{c}{8. }\\
\midrule
\endfirsthead
\caption*{\normalfont{Table \ref{tab:otm1-lmer-table-random-effects} continued}}\\
\toprule
 & \multicolumn{1}{c}{\% of variance} & \multicolumn{1}{c}{1. } & \multicolumn{1}{c}{2. } & \multicolumn{1}{c}{3. } & \multicolumn{1}{c}{4. } & \multicolumn{1}{c}{5. } & \multicolumn{1}{c}{6. } & \multicolumn{1}{c}{7. } & \multicolumn{1}{c}{8. }\\
\midrule
\endhead
Participant &  &  &  &  &  &  &  &  & \\
\ \ \ 1. Intercept & .35 & 0.76 & -0.18 & -0.38 & 0.28 &  &  &  & \\
\ \ \ 2. Response mapping & .05 &  & 0.30 & 0.01 & -0.54 &  &  &  & \\
\ \ \ 3. Learning block & .48 &  &  & 0.90 & -0.18 &  &  &  & \\
\ \ \ 4. Response mapping $\times$ Learning block & .03 &  &  &  & 0.23 &  &  &  & \\ \midrule
Stimulus &  &  &  &  &  &  &  &  & \\
\ \ \ 1. Intercept & .02 & 0.16 & 0.67 & 0.22 & -0.19 & -0.89 & 0.15 & -0.01 & 0.24\\
\ \ \ 2. Response mapping & .00 &  & 0.08 & -0.03 & -0.28 & -0.60 & 0.08 & 0.25 & -0.12\\
\ \ \ 3. Learning block & .00 &  &  & 0.05 & 0.40 & -0.60 & -0.50 & -0.94 & 0.98\\
\ \ \ 4. Valence order & .02 &  &  &  & 0.17 & 0.11 & -0.96 & -0.67 & 0.49\\
\ \ \ 5. Response mapping $\times$ Learning block & .01 &  &  &  &  & 0.09 & 0.00 & 0.36 & -0.58\\
\ \ \ 6. Response mapping $\times$ Valence order & .01 &  &  &  &  &  & 0.13 & 0.71 & -0.54\\
\ \ \ 7. Learning block $\times$ Valence order & .02 &  &  &  &  &  &  & 0.18 & -0.96\\
\ \ \ 8. Response mapping $\times$ Learning block $\times$ Valence order & .02 &  &  &  &  &  &  &  & 0.17\\
\bottomrule
\addlinespace
\insertTableNotes
\end{longtable}

}

\end{lltable}

In line with the ANOVA results, we found the expected three-way interaction between \emph{Response mapping}, \emph{Valence order}, and \emph{Learning block}, Table~\ref{tab:otm1-lmer-table-fixed-effects} and \ref{tab:otm1-lmer-table-random-effects}.
The three-way interaction prompted us to test the differences between response mapping effects in the first and second learning block for each valence order.
We found that response time differences indicated \emph{more} favorable evaluations after the first than after the second block when the behavioral information was first positive and later negative, \(\Delta M = 0.21\), 95\% CI \([0.10, 0.32]\), \(t(23.79) = 3.94\), \(p < .001\).
Vice versa, response time differences indicated \emph{less} favorable evaluations after the first than after the second block when descriptions of Bob were first negative and later positive, \(\Delta M = -0.24\), 95\% CI \([-0.34, -0.14]\), \(t(41.39) = -4.94\), \(p < .001\).
Again, these results indicate that the rating and IAT scores did not dissociate.

\hypertarget{bayesian-model-comparison}{%
\subsection{Bayesian model comparison}\label{bayesian-model-comparison}}

We implemented the unconstrained model as a hierarchical linear model that encompasses each of the other models as special cases:

\[
\begin{aligned}
\hat y_{ijk} = & \mu_i + \eta_l x_{1il} + \\
          & (\alpha_i + \tau_l x_{1il}) x_{2j} x_{3k} + \\
          & (\beta_i + \upsilon_l x_{1il}) (1 -  x_{2j}) x_{3k}
\end{aligned}
\]

The model predicts the \(i\)th participant's response to evaluation measure \(j\) in the experimental block \(k\).
Responses are predicted as a combination of participant-specific intercepts \(\mu_i\) (i.e., habitually higher or lower evaluations), a main effect of the labs \(\eta_l\), and participant-specific simple effects of learning block for rating scores (\(\alpha_i\)) and IAT score (\(\beta_i\)).
Additionally, we allowed the simple effects to be moderated by the labs (\(\tau_l\) and \(\upsilon_l\) represent the lab-specific deviations from the overall simple effects).
Participant-specific intercepts and simple effects were modeled as the sum of fixed effects and by-participant random effects, such that \(\mu_i = \mathcal{N}(\mu, \sigma_\mu)\), \(\alpha_i = \mathcal{N}(\alpha, \sigma_\alpha)\), and \(\beta_i = \mathcal{N}(\beta, \sigma_\beta)\).

The model does not include a main effect of evaluative measure because any mean differences between evaluative measures were leveled by the by-measure \(z\) standardization.
\(x_{1il}\) represents \(l\) effect coded variables that indicate which lab participant \(i\) belongs to; \(x_{2j}\) indicates the evaluative measure (1 for rating score and 0 for IAT score), such that \(\alpha_i + \tau_l\) is only relevant for rating scores and \(\beta_i + \upsilon_l\) is only relevant for IAT scores; \(x_{3k}\) is an effect coded variable that is set to 0.5 for block 1 and -0.5 for block 2.

This model allowed us to place priors on the simple effects (in units of standardized mean differences \(d\)) for each evaluative measure and implement the theoretically motivated order constraints:

\[
\begin{aligned}
\mathcal{M}_\textrm{No effect}:~ & \delta_\alpha = 0 \\ & \delta_\beta = 0 \\
\mathcal{M}_\textrm{One mind}:~ & \delta_\alpha \sim \textrm{Positive-Half-Cauchy}(r = \sqrt2/2) \\ & \delta_\beta \sim \textrm{Positive-Half-Cauchy}(r = \sqrt2/2) \\
\mathcal{M}_\textrm{Two minds}:~ & \delta_\alpha \sim \textrm{Positive-Half-Cauchy}(r = \sqrt2/2) \\ & \delta_\beta \sim \textrm{Negative-Half-Cauchy}(r = \sqrt2/2) \\
\mathcal{M}_\textrm{Any effect}:~ & \delta_\alpha \sim \textrm{Cauchy}(r = \sqrt2/2) \\ & \delta_\beta \sim \textrm{Cauchy}(r = \sqrt2/2)
\end{aligned}
\]

Additionally, we placed default multivariate Cauchy priors (\(r = \sqrt2/2\)) on lab main effects \(\eta_l\) as well as on lab effects on evaluative differences between blocks for rating scores (\(\tau_l\)) and IAT scores (\(\upsilon_l\)).

To formally assess whether the data from all labs exhibited consistent effects we added another model that enforced the order constraint of \(\mathcal{M}_\textrm{One mind}\) and \(\mathcal{M}_\textrm{Two minds}\) not only for the average block effects (\(\alpha\) and \(\beta\)) but for each lab individually (i.e., \(\alpha_l = \alpha + \tau_l\) and \(\beta_l = \beta + \upsilon_l\); \(\mathcal{M}_\textrm{One mind everywhere}\) and \(\mathcal{M}_\textrm{Two minds everywhere}\)).

For the analyses we drew 5 million samples to estimate the postrior distribtution of model parameters.
Because the draws from the posterior distribution are used to estimate the Bayes factors for model comparisons that involve order constraints (Klugkist, Laudy, \& Hoijtink, 2005), the number of draws implies upper and lower bounds on some of the reported Bayes factors.
Most notably, as a direct consequence of the number MCMC samples the \(\textrm{BF}_{\mathcal{M}_\textrm{One mind}/\mathcal{M}_\textrm{Two minds}} \in [1 / (5 \times 10^6), 5 \times 10^6]\).

\hypertarget{prior-sensitivity-analysis}{%
\subsubsection{Prior sensitivity analysis}\label{prior-sensitivity-analysis}}

Bayesian model comparison by Bayes factors are by definition sensitive to the specified prior distributions.
To ensure that our inference is not contigent on our choice of priors, we conducted prior sensitivity analyses for our key results.

Our choice of piors for the simple effects of learning block for rating scores (\(\alpha\)) and IAT score (\(\beta\)) could be viewed as either overly optimistic or pessimistic.
The prior on simple rating score effects places considerable probability mass on effects \(d\) \textless{} 0.707 although the previously reported effects were very large.
Similarly, placing the same prior on the simple effects for rating and IAT scores could be criticized because the previously reported IAT score effects were considerably smaller than those of rating scores.

We, therefore, varied the scale for the Cauchy priors on the simple effects in the ranges of \(0.50 < r_\alpha < 1.41\) and \(0.14 < r_\beta < 0.71\) for rating and IAT scores, respectively.
In light of the results from previous studies, we limited our analysis to combinations where the prior scale was larger for rating than for IAT effects.
To reduce the computational burden of this analysis, we drew 500,000 samples for each model comparison and each combination of prior scales, which implies that the Bayes factors for comparisons involving order constraints are bounded within the range of \([1 / (5 \times 10^5), 5 \times 10^5]\).
The results of the prior sensitvity analysis reassure us that our inference is robust to a wide range and combination of scales of the default Cauchy priors, see Table~\ref{tab:otm1-prior-sensitivity-replication-table}.
The Bayes factors were not affected by the scale of the priors to any meaningful degree.
This is because our data are informative enough to overwhelm the priors and because these Bayes factors depend only on the shape and location of the joint posterior distribution (Klugkist, Kato, \& Hoijtink, 2005).






\begin{table}[tbp]

\begin{center}
\begin{threeparttable}

\caption{\label{tab:otm1-prior-sensitivity-replication-table}Results of the prior sensitivity analysis for the Bayesian model comparisons of primary interest.}

\begin{tabular}{cccc}
\toprule
$r_\alpha$ & \multicolumn{1}{c}{$r_\beta$} & \multicolumn{1}{c}{$\textrm{BF}_{\mathcal{M}_\textrm{One mind}/\mathcal{M}_\textrm{Two minds}}$} & \multicolumn{1}{c}{$\textrm{BF}_{\mathcal{M}_\textrm{One mind}/\mathcal{M}_\textrm{Any effect}}$}\\
\midrule
0.50 & 0.14 & 500,001.00 & 4.00\\
0.96 & 0.14 & 500,001.00 & 4.00\\
1.41 & 0.14 & 500,001.00 & 4.00\\
0.50 & 0.42 & 500,001.00 & 4.00\\
0.96 & 0.42 & 500,001.00 & 4.00\\
1.41 & 0.42 & 500,001.00 & 4.00\\
0.96 & 0.71 & 500,001.00 & 4.00\\
1.41 & 0.71 & 500,001.00 & 4.00\\
\bottomrule
\addlinespace
\end{tabular}

\begin{tablenotes}[para]
\normalsize{\textit{Note.} The Bayes factor (BF) in favor of \(\mathcal{M}_\textrm{One mind}\) relative to \(\mathcal{M}_\textrm{Any effect}\) is bounded within the range of \([0, 4]\).
\(r_\alpha\) and \(r_\beta\) denote the scale for the Cauchy prior on the simple effects of learning block for rating scores (\(\alpha\)) and IAT scores (\(\beta\)), respectively (in units of standard deviations).}
\end{tablenotes}

\end{threeparttable}
\end{center}

\end{table}

\hypertarget{recognition-task-prior-sensitivity-analysis}{%
\subsection{Recognition task prior sensitivity analysis}\label{recognition-task-prior-sensitivity-analysis}}

To test the robustness of our inference regarding participants recognition accuracy we varied the scale \(r\) of the Cauchy prior in a wide interval of \([0.25, 1]\).
The resulting Bayes factors were \(3.89 \times 10^{6}< \textrm{BF}_{10} < 4.89 \times 10^{6}\) and thus varied by a factor of 0.80.
These results again reassure that our inference is robust to a wide range of scales of the default Cauchy prior.

\hypertarget{word-recogniton-and-iat-score-differences}{%
\subsection{Word recogniton and IAT score differences}\label{word-recogniton-and-iat-score-differences}}



\begin{figure}
\centering
\includegraphics{supplements_files/figure-latex/otm1-rydell-recognition-plot-1.pdf}
\caption{\label{fig:otm1-rydell-recognition-plot}Black-rimmed points represent condition means, error bars represent 95\% bootstrap confidence intervals based on 10,000 samples. Small points represent individual participants' accuracy. Violins represent kernel density estimates of sample distributions.}
\end{figure}

In contrast to the original results reported by Rydell et al. (2006), the recognition accuracy of briefly flashed words was above chance in this study, Figure~\ref{fig:otm1-rydell-recognition-plot}.
Memory for these words may, thus, have interfered with the associative learning process and prevented the predicted reversal of IAT score differences between learning blocks.
We, therefore, performed an exploratory regression analysis of IAT score difference and recognition accuracy of briefly flashed words.
Positive IAT score differences indicate a more favorable evaluation after the block in which we presented positive behavioral information and briefly flashed negative words.
If word recognition indeed obstructed the associative learning process, we would expect to observe a positive relationship between the recognition accuracy and IAT score differences between blocks:
When the recognition for briefly flashed words is high, IAT score differences should be positive, that is reflect the valence of the behavioral information.
We would expect to observe smaller and eventually negative IAT score differences as the recognition accuracy declines and associative learning takes over.
To account for measurement error in the recognition accuracy of briefy flashed words we fit an errors-in-variable regression model (Klauer, Draine, \& Greenwald, 1998).
Because the model assumes that predictor values are sampled from a Gaussian distribution truncated at 0, we probit-transformed the recognition accuracy.\footnote{A standard linear regression analysis yielded the same results.}




\begin{figure}
\centering
\includegraphics{supplements_files/figure-latex/otm1-prime-recognition-iat-difference-1.pdf}
\caption{\label{fig:otm1-prime-recognition-iat-difference}Scatterplot of the recognition accuracy of briefly flashed words (on probit scale) and evaluative differences in IAT scores between learning blocks in which Bob was presented with positive descriptions and those in which he was paired with negative descriptions.
The regression line and confidence band represents predictions of the errors-in-variables model (Klauer et al., 1998).}
\end{figure}

We did not detect a relationship between the recognition accuracy of briefly preseted words and IAT score differences between learning blocks, .
Moreover, the positive intercept of the regression line indicates a positive IAT score difference despite at-chance word recognition accuracy, .
Hence, even for participants who did not recognize the briefly flashed words, IAT score differences reflected the valence of the behavioral information, see Figure~\ref{fig:otm1-prime-recognition-iat-difference}.
These results provide no indication that the deviation of our findings from those reported by Rydell et al. (2006) are attributable to the above-chance recognition of briefly flashed words in this study.

\hypertarget{experiment-2}{%
\section{Experiment 2}\label{experiment-2}}

Table~\ref{tab:otm2-participant-table} summarizes the participants' demographics separately for each location of data collection.





\begin{table}[h]

\begin{center}
\begin{threeparttable}

\caption{\label{tab:otm2-participant-table}Participant demographics by location.}

\begin{tabular}{lrrrrr}
\toprule
 &  &  & \multicolumn{3}{c}{Gender} \\
\cmidrule(r){4-6}
Location & \multicolumn{1}{c}{$n$} & \multicolumn{1}{c}{Age} & \multicolumn{1}{c}{Female (\%)} & \multicolumn{1}{c}{Male (\%)} & \multicolumn{1}{c}{Non-binary (\%)}\\
\midrule
Florida & 139 & 18.40 [18, 22] & 70.50 & 28.06 & 1.44\\
Williams & 80 & 19.30 [18, 29] & 56.25 & 43.75 & 0.00\\
Hong Kong & 85 & 20.27 [18, 24] & 67.06 & 32.94 & 0.00\\
Illinois & 42 & 18.64 [18, 22] & 76.19 & 21.43 & 2.38\\
Yale & 40 & 19.15 [18, 24] & 70.00 & 27.50 & 2.50\\
\bottomrule
\addlinespace
\end{tabular}

\begin{tablenotes}[para]
\normalsize{\textit{Note.} Mean age is given with range in brackets.}
\end{tablenotes}

\end{threeparttable}
\end{center}

\end{table}

\hypertarget{prior-sensitivity-analyses}{%
\subsection{Prior sensitivity analyses}\label{prior-sensitivity-analyses}}

\hypertarget{bayesian-model-comparison-1}{%
\subsubsection{Bayesian model comparison}\label{bayesian-model-comparison-1}}

As for Experiment 1, we varied the scale for the Cauchy priors on the simple effects in the ranges of \(0.50 < r_\alpha < 1.41\) and \(0.14 < r_\beta < 0.71\) for rating and IAT scores, respectively.
The results of the prior sensitvity analysis reassure us that our inference is robust to a wide range and combination of scales of the default Cauchy priors, see Table~\ref{tab:otm2-prior-sensitivity-replication-table}.
The Bayes factors were not affected by the scale of the priors to any meaningful degree.






\begin{table}[tbp]

\begin{center}
\begin{threeparttable}

\caption{\label{tab:otm2-prior-sensitivity-replication-table}Results of the prior sensitivity analysis for the Bayesian model comparisons of primary interest.}

\begin{tabular}{cccc}
\toprule
$r_\alpha$ & \multicolumn{1}{c}{$r_\beta$} & \multicolumn{1}{c}{$\textrm{BF}_{\mathcal{M}_\textrm{One mind}/\mathcal{M}_\textrm{Two minds}}$} & \multicolumn{1}{c}{$\textrm{BF}_{\mathcal{M}_\textrm{One mind}/\mathcal{M}_\textrm{Any effect}}$}\\
\midrule
0.50 & 0.14 & 500,001.00 & 4.00\\
1.41 & 0.14 & 500,001.00 & 4.00\\
1.41 & 0.71 & 500,001.00 & 4.00\\
\bottomrule
\addlinespace
\end{tabular}

\begin{tablenotes}[para]
\normalsize{\textit{Note.} The Bayes factor (BF) in favor of \(\mathcal{M}_\textrm{One mind}\) relative to \(\mathcal{M}_\textrm{Any effect}\) is bounded within the range of \([0, 4]\).
\(r_\alpha\) and \(r_\beta\) denote the scale for the Cauchy prior on the simple effects of learning block for rating scores (\(\alpha\)) and IAT scores (\(\beta\)), respectively (in units of standard deviations).}
\end{tablenotes}

\end{threeparttable}
\end{center}

\end{table}

\hypertarget{recognition-task-prior-sensitivity-analysis-1}{%
\subsection{Recognition task prior sensitivity analysis}\label{recognition-task-prior-sensitivity-analysis-1}}



\begin{figure}
\centering
\includegraphics{supplements_files/figure-latex/otm2-rydell-recognition-plot-1.pdf}
\caption{\label{fig:otm2-rydell-recognition-plot}Black-rimmed points represent condition means, error bars represent 95\% bootstrap confidence intervals based on 10,000 samples. Small points represent individual participants' accuracy. Violins represent kernel density estimates of sample distributions.}
\end{figure}

Mirroring the original results reported by Rydell et al. (2006), the recognition accuracy of briefly flashed words was at chance in when flashed for 13 ms, Figure~\ref{fig:otm2-rydell-recognition-plot}.

To test the robustness of our inference regarding participants recognition accuracy we varied the scale \(r\) of the Cauchy prior in a wide interval of \([0.25, 1]\).
The resulting Bayes factors were \(5.13< \textrm{BF}_{01} < 19.39\) and thus varied by a factor of 3.78.
These results again reassure that our inference is robust to a wide range of scales of the default Cauchy prior.

\newpage

\hypertarget{references}{%
\section{References}\label{references}}

\begingroup
\setlength{\parskip}{0pt}
\setlength{\parindent}{-0.5in}
\setlength{\leftskip}{0.5in}

\hypertarget{refs}{}
\begin{CSLReferences}{1}{0}
\leavevmode\vadjust pre{\hypertarget{ref-Klauer_Draine_Greenwald_1998}{}}%
Klauer, K. C., Draine, S. C., \& Greenwald, A. G. (1998). An unbiased errors-in-variables approach to detecting unconscious cognition. \emph{British Journal of Mathematical and Statistical Psychology}, \emph{51}(2), 253--267. \url{https://doi.org/10.1111/j.2044-8317.1998.tb00680.x}

\leavevmode\vadjust pre{\hypertarget{ref-klugkist_bayesian_2005}{}}%
Klugkist, I., Kato, B., \& Hoijtink, H. (2005). Bayesian model selection using encompassing priors. \emph{Statistica Neerlandica}, \emph{59}(1), 57--69. \url{https://doi.org/10.1111/j.1467-9574.2005.00279.x}

\leavevmode\vadjust pre{\hypertarget{ref-klugkist_inequality_2005}{}}%
Klugkist, I., Laudy, O., \& Hoijtink, H. (2005). Inequality constrained analysis of variance: A bayesian approach. \emph{Psychological Methods}, \emph{10}(4), 477--493. \url{https://doi.org/10.1037/1082-989x.10.4.477}

\leavevmode\vadjust pre{\hypertarget{ref-rydell_two_2006}{}}%
Rydell, R. J., McConnell, A. R., Mackie, D. M., \& Strain, L. M. (2006). Of {Two Minds}: {Forming} and {Changing Valence}-{Inconsistent Implicit} and {Explicit Attitudes}. \emph{Psychological Science}, \emph{17}(11), 954--958. \url{https://doi.org/10.1111/j.1467-9280.2006.01811.x}

\leavevmode\vadjust pre{\hypertarget{ref-wolsiefer_modeling_2017}{}}%
Wolsiefer, K., Westfall, J., \& Judd, C. M. (2017). Modeling stimulus variation in three common implicit attitude tasks. \emph{Behavior Research Methods}, \emph{49}(4), 1193--1209. \url{https://doi.org/10.3758/s13428-016-0779-0}

\end{CSLReferences}

\endgroup


\end{document}
